%% Overleaf: The \marks command has been renamed to 
%% \Qmarks to enable correct compilation on recent LaTeX
%% distributions (30 September 2016)
\documentclass[11pt,qn]{examform}

\usepackage{graphicx}

\def\subjecttitle{The \LaTeX\ Typesetting System / LKV101}
\def\semester{2002/3(Sem 2)}
\def\timeallowed{2.5h}
\def\examiner{Ling Keck Voon}
\def\examiners{Ling KV \& Co.}
\def\moderator{Name of Moderator here}
\def\subjectmoderator{Name of Subject Moderator here}

\begin{document}

\begin{question}{1}
A \LaTeX\ template for typesetting questions. It is implemented
using the {\tt fancyhdr} style and {\tt longtable} environment
which comes standard with the \LaTeX\ tools package. This is
version 0.00 and it is very crude. An customised environment
called {\em question} is provided. You just need to say \\

\begin{verbatim}
\documentclass[11pt,qn]{examform}

\usepackage{graphicx}

\def\subjecttitle{The \LaTeX\ Typesetting System / LKV101}
\def\semester{2002/3(Sem 2)}
\def\timeallowed{2.5h}
\def\examiner{Ling Keck Voon}
\def\examiners{Ling KV \& Co.}
\def\moderator{Name of Moderator here}
\def\subjectmoderator{Name of Subject Moderator here}

\begin{document}

\begin{question}{1}
  Type your question 1 here....
\end{question}
\begin{question}{2}
  Type your question 2 here...
\end{question}
\end{document}
\end{verbatim} \\

\newpage

At the moment, you need to manually adjust the vertical white
spaces so that chunks of text/graphics fit well into a page. This
is how you do it: \\

\begin{verbatim}
\vspace{10mm} \\
\end{verbatim}

Comments and tricks welcome.\\


Marks can be typeset by the command \verb^\Qmarks{num}^ where {\tt
num} a number. \\

If you could typeset this equation, you would be awarded 10 marks.
\[ g(s) = \frac{k}{\tau s +1}e^{-\theta s} \]
\Qmarks{10}\\

Five marks for including a figure.
\begin{center}
\includegraphics[width=0.5\textwidth]{fig1.pdf}
\end{center}
\Qmarks{5}\\

\end{question}
%%%%%%%%

\end{document}
